\section{Radical-step fibres, global prime packets, and the same-pilot type boundary}
\label{sec:global-packets-same-pilot}

This section continues the positive and negative searches simultaneously.
Failure of a finite search, absence of a known theorem, and present technical
difficulty do not retire a route.  A counterexample below is used only against
the fully quantified statement that it satisfies; the surrounding route is
kept whenever a stronger or differently typed mechanism survives.

\subsection{Replacing the product step by the seed radical}

Let $a+b=c$ be primitive, put $P=abc$ and $R=\rad(P)$, and let $Q$ be any
positive multiple of $R$.  For $h,k\geq1$ define
\begin{equation}\label{eq:gpsp-radical-shear}
 U=1+Qh,\qquad V=1+Q(h+ck),\qquad W=1+Q(h+bk).
\end{equation}

\begin{theorem}[Radical-step primitive shear]
\label{thm:gpsp-radical-shear}
If $\gcd(U,k)=1$, then $(aU,bV,cW)$ is a primitive $abc$ triple.  Each of
the three pair projections $(U,V)$, $(U,W)$, and $(V,W)$ determines $(h,k)$.
If $K\geq2$ and
\[
 M_Q=\left\lfloor\frac{c^{K-2}}{4Q}\right\rfloor,
\]
then every parameter pair $1\leq h,k\leq M_Q$ has height below $c^K$.
For sufficiently large $M_Q$, at least $M_Q^2/8$ of these pairs are
admissible.
\end{theorem}

\begin{proof}
Every prime dividing a seed coordinate divides $Q$, while all three
cofactors are congruent to one modulo $Q$.  Thus every cofactor avoids the
entire seed support.  The identities
\[
 V-U=cQk,\qquad W-U=bQk,\qquad V-W=aQk
\]
show that a common divisor of $U$ and $V$, or of $U$ and $W$, must divide
$k$ after the factors $Q$ and the appropriate seed coordinate have been
removed.  The hypothesis $\gcd(U,k)=1$ excludes both possibilities.  A
common divisor of $V$ and $W$ divides $aQk$; it avoids $aQ$, and
$V\equiv U\pmod{k}$, so it also avoids $k$.  Hence $U,V,W$ are pairwise
coprime.  Combining this with support avoidance and the pairwise
coprimality of $a,b,c$ proves primitivity.  The identity
$aU+bV=cW$ follows by expansion.

Any one of the three displayed differences recovers $k$ from the indicated
pair.  If the pair contains $U$, then $U=1+Qh$ recovers $h$; from $(V,W)$ one
instead uses $h=(V-1)/Q-ck$.  This proves the three injectivity statements.
Since $b+1\leq c$,
\[
 W\leq1+QcM_Q\leq1+\frac{c^{K-1}}4,
\]
and therefore $cW<c^K$ for $c\geq2$ and $K\geq2$.  Finally, the elementary
sieve used in Lemma~\ref{lem:bp-admissible-count} applies verbatim with $Q$
in place of $P$: the bad pairs are covered by primes dividing both $k$ and
$1+Qh$.  It leaves at least $M_Q^2/8$ pairs once $M_Q$ is sufficiently
large.
\end{proof}

The minimal step $Q=R$ is therefore optimal among these multiples for the
raw box.  At $K=8$ its elementary lower count has scale
\begin{equation}\label{eq:gpsp-raw-rad-step}
                    \frac{c^{12}}{R^2},
\end{equation}
instead of $c^{12}/P^2$.  This is a genuine improvement when the seed has
large prime powers.  It does not itself count low-radical outputs.  Indeed,
the pair-projection upper argument is intrinsic to the output factors and
remains
\begin{equation}\label{eq:gpsp-exception-upper}
 \#\mathcal E_Q(c^8)\ll_{\varepsilon}R^{-2/3}c^{4+\varepsilon}
\end{equation}
on the $\mu=3/4$ locus.

There is no gain from averaging the same construction over multiples of
$R$.  If $Q=sR$, then the $Q$-point at $(h,k)$ is exactly the $R$-point at
$(sh,sk)$.  Moreover $1+Qh\equiv1\pmod{s}$, so admissibility at $(h,k)$
implies admissibility at $(sh,sk)$.  Thus the larger-step fibres are
subfibres of the minimal one after rescaling their boxes.

The often stated ``matching lower bound'' should also be interpreted with
care.  Assume the uniform upper estimate \eqref{eq:gpsp-exception-upper}.
For fixed $\lambda<1$ and $\eta>0$, the assertion that there are constants
$A>0,c_0$ such that every seed with $R<c^\lambda$ and $c\geq c_0$ satisfies
\begin{equation}\label{eq:gpsp-eventual-lower}
 \#\mathcal E_R(c^8)\geq A R^{-2/3}c^{4+\eta}
\end{equation}
is equivalent to boundedness of the heights of the seeds in that
subcritical locus.  The forward implication compares
\eqref{eq:gpsp-eventual-lower} with \eqref{eq:gpsp-exception-upper} at
$\varepsilon=\eta/2$.  Conversely, after a height bound has already been
proved, one may choose $c_0$ above it and the lower assertion is vacuous.
Consequently \eqref{eq:gpsp-eventual-lower} is an exact closing gate, but it
is not by itself an explanatory distribution theorem.

A fixed congruence template cannot provide that explanation.  Suppose
$d_1,d_2,d_3$ are pairwise coprime, $D=d_1d_2d_3$, and
$\gcd(D,Qabc)=1$.  The conditions
\[
                   d_1\mid U,\qquad d_2\mid V,
                   \qquad d_3\mid W
\]
have exactly $D$ solutions $(h,k)$ modulo $D$, hence density $D^{-1}$.
The deterministic lower certificate for repeated-prime excess supplied
solely by these divisibilities is
\begin{equation}\label{eq:gpsp-crt-excess}
             \Delta=\prod_{i=1}^3\frac{d_i}{\rad(d_i)}\leq D.
\end{equation}
On the upper half of the $K=8$ radical-step box, forcing the target radical
inequality requires $\Delta\gg Rc^{14}$.  It therefore also requires
$D\gg Rc^{14}$, while the periodic main term is at most
\[
       \frac{M_R^2}{D}\ll \frac{c^{12}/R^2}{Rc^{14}}
        =\frac1{c^2R^3}.
\]
This rules out obtaining the matching lower bound from the periodic main term
of a single template alone.  It leaves
large unions of templates, boundary terms, and accidental high valuations
open.

There are also exact finite warnings against replacing excess by a visual
``square cofactor'' condition.  For the subcritical seed $(1,8,9)$, $R=6$,
the choices
\[
 (Q,h,k)=(6,840,60)\quad\hbox{and}\quad(72,70,5)
\]
both give
\[
 (U,V,W)=(5041,8281,7921)=(71^2,91^2,89^2).
\]
Nevertheless the output height is $71289$, its radical is $3450174$, and
the desired fourth-power versus third-power radical inequality fails.  A
second example is
\[
 (Q,h,k)=(72,1160,798),\qquad
 (U,V,W)=(289^2,775^2,737^2),
\]
and it fails the same target.  These are counterexamples to the assertion
that three square cofactors automatically give an exceptional output.  They
do not refute the eventual arithmetic lower gate.

\subsection{Exact order stagnation in the balancing orbit}

Let
\[
 u_0=0,\quad u_1=1,\quad u_{n+2}=6u_{n+1}-u_n,
\]
and put $\alpha=3+2\sqrt2$ and $\gamma=\alpha^2=17+12\sqrt2$.  Then
\begin{equation}\label{eq:gpsp-geometric-pell}
 u_n=\alpha^{1-n}\frac{\gamma^n-1}{\gamma-1}.
\end{equation}
For an odd rational prime $q$, let $z(q)$ be its rank of apparition and put
$e(q)=v_q(u_{z(q)})$.

\begin{theorem}[Order-stagnation tower]
\label{thm:gpsp-order-stagnation}
For every odd prime $q$ and $k\geq1$,
\begin{equation}\label{eq:gpsp-order-stagnation}
 z(q^k)=z(q)q^{\max(0,k-e(q))}.
\end{equation}
In particular, depth two is exactly one failure of order growth and depth
three is exactly two consecutive failures.
\end{theorem}

\begin{proof}
In $K=\mathbb Q(\sqrt2)$, every odd $q$ is unramified and
$N(\gamma-1)=-32$.  At every prime ideal $\mathfrak q\mid q$,
\eqref{eq:gpsp-geometric-pell} identifies $z(q)$ with the residual order of
$\gamma$ and $e(q)$ with
$v_{\mathfrak q}(\gamma^{z(q)}-1)$.  If $y$ is a local unit with
$v(y-1)=e\geq1$, the odd-prime binomial expansion gives
\[
                         v(y^m-1)=e+v_q(m).
\]
Thus the least multiplier needed to raise the congruence from level one to
level $k$ is $q^{\max(0,k-e)}$.  Substitution gives
\eqref{eq:gpsp-order-stagnation}.
\end{proof}

At an odd prime index $\ell$, write
$(1+\sqrt2)^\ell=A_\ell+B_\ell\sqrt2$.  The previous packet theorem says
that a squarefull $u_\ell=A_\ell B_\ell$ forces four distinct rank-$\ell$
repeated primes, two in each coprime channel, with a depth-three prime in
each channel.  There is also a first-order relation among all packet primes.
If
\[
 A_\ell=1+2\ell a,\qquad
 \left(\frac2\ell\right)B_\ell=1+2\ell b,
\]
then the Pell norm gives the exact identity
\begin{equation}\label{eq:gpsp-channel-ledger}
                    a-2b+\ell(a^2-2b^2)=0.
\end{equation}
Expanding the prime factorizations modulo $4\ell^2$ turns its reduction
modulo $\ell$ into a coupled congruence between the quotient coordinates in
the two channels.  This is stronger than independent residue conditions,
although the free quotient variables prevent a contradiction at present.

The latest number-field input yields a useful global alternative.  Fellini
and Murty state that finiteness of base-$x$ super-Wieferich prime ideals
forces infinitely many non-Wieferich prime ideals~\cite{FelliniMurty2026}.
An adversarial audit of their Section~8 architecture gives a sharper
prime-order consequence, but it also requires two local proof repairs and one
textual correction.  At a selected prime index $\ell$, a prime ideal
$\mathfrak p\mid\Phi_\ell(x)$ has residual order dividing $\ell$.  Order one
would imply both $x\equiv1$ and $\ell\equiv0$ modulo $\mathfrak p$; the
defining exclusion of primes above $\ell$ which divide $(x-1)$ rules this
out.  Thus every such factor has order exactly $\ell$, and its exponent in
$\Phi_\ell(x)$ is its first-occurrence valuation.  Excluding the finitely
many super-Wieferich orders leaves exponents one or two.  If only finitely
many prime-order exponent-one ideals existed, the Section~8 localization,
finite square-class decomposition, and Siegel argument would again place
infinitely many distinct points on a finite family of genus-one curves.

This streamlined prime-index proof deliberately avoids the printed second
case of Lemma~6.4.  The displayed equality there is generally false: with
$x=2,p=3,f=2,i=1$, it identifies $\Phi_6(2)=3$ with
$(2^6-1)/(2^2-1)=21$.  Its binomial proof also treats
$\binom pp=1$ as if it were divisible by $p$; separating the last term
repairs the valuation estimate.  A contradictory sentence on printed
p.~226 is avoided by allowing the localized square generator to be a unit.
These counterexamples invalidate intermediate printed steps.  They do not
invalidate the repaired valuation conclusion or the resulting Siegel
argument.

Specializing this proof extraction to $x=\gamma$ and descending prime ideals
to rational primes gives the unconditional alternative.  After discarding
the finitely many prime ideals above $2$, the local dictionary applies to odd
rational primes:
\begin{equation}\label{eq:gpsp-global-alternative}
\begin{split}
&\text{either infinitely many odd primes $q$ satisfy $e(q)\geq3$,}\cr
&\text{or infinitely many odd prime indices $\ell$ admit
  $q\parallel u_\ell$ with $z(q)=\ell$.}
\end{split}
\end{equation}
This is a repaired refinement of the cited proof architecture, not a
quotation of its theorem statement.  In the first branch it gives neither a common rank nor opposite
channels; in the second it gives infinitely many nonsquarefull prime-index
terms but not all prime indices.  It therefore does not settle the forced
four-prime packet.

Three natural shortcuts have exact counterexamples.  At index seven,
\[
 B_7=13^2,\qquad 13^2\parallel u_7,\qquad z(13)=7,
 \qquad13=2\cdot7-1.
\]
Moreover
\[
 F_7(X)=X^6-5X^4+6X^2-1,\quad
 F_7(6)=40391=13^2\cdot239,\quad F_7'(6)\equiv2\pmod{13}.
\]
Thus a primitive divisor need not be simple, equality in the minimal channel
bound does not force simplicity, and a nonsingular polynomial root modulo
$q$ does not force the fixed integer value to avoid $q^2$.  The opposite
channel has the further exact example
\[
 1546463=2\cdot773231+1,\quad z(1546463)=773231,\quad
 1546463^2\parallel A_{773231},
\]
with both displayed rank parameters prime.  These examples close only the
three stated implications.

Two independent exact programs exhaust all $5{,}761{,}454$ odd primes
$q\leq10^8$.  The only repeated primes are $13,31,1546463$, all of exact
depth two; there is no depth-three hit.  Hence both depth-three primes in a
hypothetical squarefull prime-index packet exceed $10^8$.  This is a finite
certificate and is never used as an infinitude conclusion.

\subsection{Recursive prime-square lifts on the Danilov progression}

On the surviving Danilov progression, write $t=T+Qr$ and maintain
\begin{equation}\label{eq:gpsp-dan-state}
                         3T+1=hQ,\qquad h\in\{1,2\}.
\end{equation}
The quadratic-unit orbit gives the exact Fibonacci factorization
\begin{equation}\label{eq:gpsp-dan-fib}
                         L_T=5F_{10hQ}F_{10hQ-5}.
\end{equation}
If $p\parallel F_{10Q}$, $p\nmid3375Q$, then the packet modulo $p^2$ is
nondegenerate.  Squarefullness forces the unique residue
\begin{equation}\label{eq:gpsp-dan-root}
                         h+3r\equiv0\pmod p.
\end{equation}
For its representative $0\leq\rho<p$, divisibility and the inequalities
$0<h+3\rho<3p$ give a unique $h'\in\{1,2\}$ with
$h+3\rho=h'p$.  The update
\[
                 T'=T+Q\rho,\qquad Q'=Qp
\]
then preserves $3T'+1=h'Q'$.

Carmichael's theorem in Yabuta's form supplies a primitive divisor of
$F_{10Q}$ outside its explicit small exceptions~\cite{Yabuta2001}.  Rank of
apparition makes such a divisor fresh relative to $Q$.  The theorem does not
control its first valuation.  This distinction yields the exact conditional
closure.

\begin{theorem}[Conditional simple-primitive closure]
\label{thm:gpsp-dan-conditional}
Assume that at every adaptively generated state, $F_{10Q}$ has a primitive
prime divisor $p$ with $p\parallel F_{10Q}$.  Then no index in the initial
Danilov survivor progression has squarefull remainder.
\end{theorem}

\begin{proof}
The simple divisor supplies a nondegenerate fresh packet and
\eqref{eq:gpsp-dan-root} nests the next progression while preserving
\eqref{eq:gpsp-dan-state}.  Repetition constructs arbitrarily many distinct
fresh primes.  If one fixed index $t$ survived as squarefull, every chosen
prime would divide the one fixed nonzero integer $L_t$.  A finite integer has
only finitely many distinct prime divisors, a contradiction.
\end{proof}

The valuation-one hypothesis is essential to this proof and is not a
consequence of the cited primitive-divisor theorem.  Nor does one abstract
packet automatically generate another.  In $\mathbb Z[\sqrt5]$, take
\[
 \eta=(9+4\sqrt5)^8,\qquad \alpha_0=19+\sqrt5,\qquad
 L_r=K_r=2\operatorname{Re}(\alpha_0\eta^r)+11.
\]
Modulo $49$, $\eta=1+35\sqrt5$ and
\begin{equation}\label{eq:gpsp-dan-countermodel}
                              L_r\equiv7r\pmod{49}.
\end{equation}
At $(T,Q,p)=(0,1,7)$ the slope is nonzero, the forced root is $r=0$, and
$K_0=49$ is squarefull.  The updated modulus is $7$, while every prime
divisor of $L_0=49$ already divides it.  Hence no fresh successor exists.
This satisfies the full abstract local-packet mechanism, but its initial norm
is not the Danilov norm.  It refutes automatic abstract continuation and does
not refute the Danilov-specific simple-primitive route.

An exact adaptive computation constructs $626$ distinct packets.  They occur
in thirteen nonempty batches.  The final state has a $4398$-digit modulus and
exactly $638$ distinct prime factors; no next eligible packet occurs below
$10^8$.  Every selected prime has exact valuation one, and every CRT update
is independently replayed.  This certifies the finite recursion, while the
endpoint remains open above its bound.

\subsection{The simple-primitive gate and Wall--Sun--Sun support}

The missing valuation-one premise can be characterized exactly for the
Fibonacci indices arising above.  If $n>5$, $5\mid n$, and $p$ is a primitive
divisor of $F_n$, then the rank of apparition satisfies
\[
 z(p)=n\mid p-\left(\frac5p\right).
\]
Reduction modulo five and quadratic reciprocity exclude the inert sign, so
\begin{equation}\label{eq:gpsp-split-primitive}
                         p\equiv1\pmod n.
\end{equation}
Sanna's exact valuation formula then gives
\begin{equation}\label{eq:gpsp-wss-transfer}
                  v_p(F_n)=v_p(F_{p-1}).
\end{equation}
Thus a primitive divisor is repeated precisely when it is a
Fibonacci--Wieferich, or Wall--Sun--Sun, prime~\cite{Sanna2016,BragmanRowland2025}.
No such prime is presently known; that empirical fact is not used as a
nonexistence theorem.

Let $Q_*$ be the final $4398$-digit squarefree modulus in the certified
recursive chain and put $n_*=10Q_*$.  Its complete factorization consists of
exactly $638$ primes, each below $10^8$, and $\gcd(Q_*,30)=1$.  Let $C_n$
denote the Fibonacci cyclotomic factor.  For $n>2$ divisible by $10$,
Yabuta's correction analysis says that $C_n$ has at most one nonprimitive
prime factor; if present, it is the greatest prime factor of $n$ and occurs
once.  At $n_*$ this correction cannot occur.  Indeed,
squarefreeness would give $n_*=qz(q)$, whence
$n_*\le q(q+1)\le10^8(10^8+1)$, contrary to the number of digits.

\begin{theorem}[Exact surviving simple-primitive failure]
\label{thm:gpsp-wss-packet}
If $F_{n_*}$ has no simple primitive divisor, then $C_{n_*}$ is powerful and
every prime $p\mid C_{n_*}$ is Wall--Sun--Sun and satisfies
$p\equiv1\pmod {n_*}$.  At least one such prime satisfies
\[
                              p\ge41n_*+1.
\]
\end{theorem}

\begin{proof}
The correction just excluded shows that every factor of $C_{n_*}$ is
primitive with its full $F_{n_*}$ multiplicity.  Equations
\eqref{eq:gpsp-split-primitive} and \eqref{eq:gpsp-wss-transfer} therefore
give the powerful and Wall--Sun--Sun assertions.  Hong's explicit large
primitive-divisor theorem with $\kappa=40$ applies because
$\log n_*>4398\log10>10036$~\cite{Hong2025}.  It produces a primitive prime
outside $jn_*\pm1$ for $1\le j\le40$.  The split congruence leaves only the
plus sign, so its coefficient is at least $41$.
\end{proof}

This theorem sharply locates the surviving bad case but does not contradict
it.  A broad real-Lucas shortcut is false under its full standard
hypotheses: for
\[
 U_{m+2}=2U_{m+1}+3U_m,
 \qquad U_{10}=14762=2\cdot11^2\cdot61,
\]
the prime $11$ is the unique primitive divisor at index ten and occurs
exactly twice.  The roots are $3,-1$, so the sequence is real,
nondegenerate and has coprime parameters.  This counterexample closes only
the sequence-uniform transfer; it does not satisfy the Fibonacci-specific
hypothesis.  An exact search over the $252$ squarefree $Q\le1000$ with
$\gcd(Q,30)=1$ finds simple primitive certificates for $207$ cases below
$5\cdot10^7$.  The remaining $45$ cases are unresolved by that finite bound,
not counterexamples.

\subsection{The current LANA specification and the pointed same-pilot gate}

The public Project LANA Lean repository was audited at the pinned commit
\texttt{ddaddc2}~\cite{LANAlean}.
Its declared Corollary~3.12 variant shares initial theta data between a
$q$-pilot scalar and an independently supplied right-hand-side container.
At that commit the right-hand-side record has a prior consistency defect.

\begin{theorem}[Empty-set obstruction in the pinned low-resolution signature]
\label{thm:gpsp-lana-empty}
For every initial theta datum $D$, the pinned type \texttt{RHSData D} is
uninhabited.
\end{theorem}

\begin{proof}
The admissible prime satisfies $\ell\geq5$, and the required standard
procession therefore has positive length.  Choose one capsule index and the
rational prime $p=2$.  The record asserts that the finite sum of the weights
over the $p$-fibre is one, so that fibre cannot be empty.  A constant function
from the capsule labels to a point of the fibre supplies a component $c$.

For this component the record further asserts, for every set $U$, that
\begin{equation}\label{eq:gpsp-lana-scaling}
 \operatorname{vol}((x\mapsto px)^{-1}U)
   =\operatorname{vol}(U)+\log p.
\end{equation}
Take $U=\varnothing$.  Its inverse image is empty, so cancellation in
\eqref{eq:gpsp-lana-scaling} gives $\log2=0$, contrary to $\log2>0$.
\end{proof}

Thus the assembled variant-data type, which contains an \texttt{RHSData}
field, is also empty.  A universal theorem over that type would be vacuous,
and there can be no counterexample satisfying all its fields.  This result
rules out the pinned low-resolution signature as a satisfiable specification.
It says nothing against a corrected Haar-volume interface, the published IUT
argument, or $abc$.  A direct repair restricts
\eqref{eq:gpsp-lana-scaling} to nonempty measurable regions of finite nonzero
volume and proves stability of that domain under preimages; alternatively one
may use an extended-real logarithmic codomain with explicit zero-volume
semantics.

After this repair, the main typing gap remains.  The two sides currently
share the initial datum but have no field realizing the $q$-pilot as a region
in the right-hand container, no identity for that region's volume, no bridge
between the two weight systems, no multiradial output map, no IPL/SHE/APT or
Ind3 data, no determinant normalization, and no volume monotonicity theorem.
Sharing a parameter supplies no order relation between two independently
chosen real outputs.

The July 2026 LANA report isolates a stronger target: for a suitable output
choice $S$, prove equality of maps
\[
        \eta^q=\eta_S^{\mathrm{anab}}:
        \mathbb R^{\mathrm{val}}\longrightarrow\mathbb R^{\mathrm{ss}}.
\]
The report states that this equality is not currently proved and does not
declare it manifestly false~\cite{LANA}.  The numerical conclusion needs
less.

\begin{proposition}[Minimal pointed same-pilot certificate]
\label{prop:gpsp-pointed-hit}
Let $p_q$ be the canonical $q$-pilot point, let $\nu$ be the correctly
normalized volume coordinate, and put
$q_*=-|\log q|$, $T=-|\log\Theta|$.  Suppose a globally synchronized,
source-permitted output $S$ satisfies
\[
 \nu(\eta^q(p_q))=q_*,\qquad
 \eta^q(p_q)=\eta_S^{\mathrm{anab}}(p_q),\qquad
 \nu(\eta_S^{\mathrm{anab}}(p_q))\leq T.
\]
Then $q_*\leq T$.
\end{proposition}

\begin{proof}
Apply $\nu$ to the pointed equality and substitute the two inequalities.
\end{proof}

The proposition is elementary after its hypotheses are available.  Its
content is the construction of the middle equality before taking the volume
quotient: the diagram must start with the fixed $q$-pilot arithmetic line,
follow the theta link, IPL/SHE/APT construction, the required
indeterminacies, log-Kummer correction and determinant normalization, and
return to the same right-hand arithmetic holomorphic structure with compatible
local and global metrics.  Defining the quotient equality by equality of the
desired scalar values would be circular.  A source-defined witness in
$\Phi(P_q)$ also suffices, but membership in that class already contains the
lower volume inequality and must be constructed independently.

The pinned specification failure is therefore a narrow no-go theorem, while
the pointed same-pilot construction remains an active positive route.  The
public report's full map equality is a stronger active target.  None of these
results yields an unconditional term of type \texttt{ABCConjecture} or its
negation.
